\begin{practice}
    \subsection{Luyện tập}

    \begin{questions}[series=SBT]
        \item
        \begin{questions}
            \item Vẽ tam giác $ABC$ vuông tại $A$, $AB=3$ cm, $AC=4$ cm. Tính $BC$, $\sin B$, $\cos B$.

            \answerline
            \answerline

            \item Vẽ tam giác $MNP$ vuông tại $M$, $MN=6$ cm, $MP=8$ cm. Hỏi hai tam giác $ABC$, $MNP$ có đồng dạng không? Tính $\sin N$, $\cos N$.

            \answerline
            \answerline
        \end{questions}

        \item
        \begin{questions}
            \item Chứng minh rằng với mọi góc nhọn $\alpha<45^\circ$, ta có
            \[
                \sin(45^\circ-\alpha)=\cos(45^\circ+\alpha),\qquad
                \cos(45^\circ-\alpha)=\sin(45^\circ+\alpha).
            \]

            \answerline
            \answerline

            \item Không dùng MTCT, tính
            \[
                \sin25^\circ+\sin35^\circ+\sin45^\circ-\cos45^\circ-\cos55^\circ-\cos65^\circ.
            \]

            \answerline
        \end{questions}

        \item Khi góc $\alpha$ lần lượt bằng $10^\circ$, $20^\circ$, $30^\circ$, $40^\circ$, hãy dùng MTCT tính $\sin\alpha$ trong mỗi trường hợp (làm tròn đến chữ số thập phân thứ ba).

        \answerline

        \item Hãy dùng MTCT, tìm số đo của góc nhọn $\alpha$ (làm tròn đến độ) trong mỗi trường hợp
        \begin{questions}
            \item Khi $\sin\alpha$ lần lượt bằng $\dfrac{1}{4}$, $\dfrac{1}{3}$, $\dfrac{1}{2}$, $\dfrac{2}{3}$;
            \item Khi $\cos\alpha$ lần lượt bằng $\dfrac{1}{4}$, $\dfrac{1}{3}$, $\dfrac{1}{2}$, $\dfrac{2}{3}$.
        \end{questions}

        \answerline
        \answerline

        \item Biết rằng với mỗi góc nhọn $\alpha$, ta có $\sin^2\alpha+\cos^2\alpha=1$, không dùng MTCT, hãy tính $\sin^225^\circ+\sin^235^\circ+\sin^245^\circ+\sin^255^\circ+\sin^265^\circ$.

        \answerline
        \answerline

        \item Một tam giác vuông có hai cạnh góc vuông đo được $5$ cm, $12$ cm. Hỏi sin của góc nhọn nhỏ nhất của tam giác đó bằng bao nhiêu?

        \answerline

        \item Xét tam giác $ABC$ vuông tại $B$, có $\widehat A=30^\circ$. Tia $Bt$ sao cho $\widehat{CBt}=30^\circ$ cắt tia $AC$ ở $D$, $D$ nằm giữa $A$ và $C$. Chứng minh rằng khoảng cách từ $D$ đến đường thẳng $BC$ bằng $\dfrac{AB}{4}$.

        \answerline
        \answerline

        \item Vẽ góc $\alpha$ trong mỗi trường hợp:
        \begin{questions}
            \item $\cos\alpha=0{,}4$;
            \item $\tan\alpha=\dfrac{2}{3}$;
            \item $\cot\alpha=\dfrac{3}{4}$.
        \end{questions}

        \answerline
        \answerline

        \item Một cái thang dài $3{,}2$ m đặt tựa bức tường, đầu thang đặt đến độ cao $3$ m thì thang tạo với mặt đất góc $\alpha$ xấp xỉ bằng bao nhiêu độ (H.4.6)?

        % TODO(HINH-NGUON): bo sung asset/crop dung Hinh 4.6 SBT trang 45 (cai thang tua tuong, nhan 3,2; 3; alpha); khong redraw xap xi.
        \answerline

        \item Một cái diều có dây diều dài $8$ m, khi dây diều căng thì diều bay ở độ cao $6$ m. Hỏi khi đó dây diều tạo với phương ngang của mặt đất góc nhọn $\alpha$ xấp xỉ bằng bao nhiêu độ (H.4.7)?

        % TODO(HINH-NGUON): bo sung asset/crop dung Hinh 4.7 SBT trang 46 (day dieu 8 m, do cao 6 m, goc alpha); khong redraw xap xi.
        \answerline

        \item Chứng minh tam giác vuông có một góc nhọn có tang bằng $1$ là tam giác vuông cân.

        \answerline
        \answerline

        \item
        \begin{questions}
            \item Tính các góc của tam giác vuông có một góc nhọn có tang bằng $\dfrac{\sqrt3}{3}$.

            \answerline

            \item Một hình chữ nhật có kích thước $3$ và $\sqrt3$. Tính các góc tạo bởi đường chéo và cạnh của hình chữ nhật đó.

            \answerline
        \end{questions}

        \item Dùng MTCT, hãy tìm tang và côtang của góc nhọn $\alpha$ khi $\alpha$ lần lượt bằng $10^\circ$, $20^\circ$, $30^\circ$, $40^\circ$ (làm tròn đến chữ số thập phân thứ ba).

        \answerline

        \item Tính tang, côtang của góc kề đáy của tam giác cân biết cạnh đáy dài $8$ cm, đường cao ứng với đáy dài $5$ cm.

        \answerline
        \answerline

        \item Dùng định nghĩa tỉ số lượng giác $\sin\alpha$, $\cos\alpha$, $\tan\alpha$, $\cot\alpha$ hãy chứng minh rằng:
        \begin{questions}
            \item $\tan\alpha=\dfrac{\sin\alpha}{\cos\alpha}$; $\cot\alpha=\dfrac{\cos\alpha}{\sin\alpha}$;
            \item $1+\tan^2\alpha=\dfrac{1}{\cos^2\alpha}$.
        \end{questions}

        \answerline
        \answerline

        \item Cho góc $\alpha$ có $\tan\alpha=\dfrac{3}{4}$. Tính $\sin\alpha$, $\cos\alpha$.

        \answerline

        \item Với $\alpha<\beta<90^\circ$, chứng minh rằng:
        \begin{questions}
            \item $\cos\alpha>\cos\beta$ (\textit{HD. Sử dụng Ví dụ 5 và bài 4.15});
            \item $\sin\alpha<\sin\beta$ (\textit{HD. Sử dụng công thức $\sin^2\alpha+\cos^2\alpha=1$}).
        \end{questions}

        \answerline
        \answerline
    \end{questions}
\end{practice}
